%% ma: L12-B08-0142
%% lop: 12
%% chuong: 2
%% bai: 8
%% dang: 12.8.48
%% loai: TN4
%% muc: VD
%% dap_an: "A"
%% nguon: "CĐ2.3.pdf giáo viên gửi 27/10/2026 (M07, Lê Minh Khôi – Nguyễn Anh Khôi, Chương II Chủ đề 3 Biểu thức tọa độ của các phép toán vectơ), Câu 58"
%% kien_thuc: "Bài 8 (điều kiện vuông góc, ba điểm thẳng hàng)"
%% trang_thai: chinh-thuc
%% ghi_chu: "Giáo viên đã duyệt 27/10/2026"
\begin{ex}
Trong không gian $Oxyz$, cho tam giác $DBK$ với tọa độ ba đỉnh là $D(2;5;3)$, $B(6;12;7)$ và $K(4;11;27)$. Gọi $H$ là hình chiếu của $K$ lên cạnh $DB$. Biết tọa độ của điểm $H$ là $H(a;b;c)$. Tính giá trị của biểu thức $T=162\cdot(a+b+c)$.
\choice{\True $T=6000$}{$T=2000$}{$T=3000$}{$T=4000$}
\loigiai{$\overrightarrow{DB}=(4;7;4)$, $DB^2=81$, $\overrightarrow{DK}=(2;6;24)$, $t=\dfrac{\overrightarrow{DK}\cdot\overrightarrow{DB}}{81}=\dfrac{146}{81}$.

$a+b+c=10+15t=10+\dfrac{730}{27}=\dfrac{1000}{27}$, $T=6000$.}
\end{ex}
